\documentclass[14pt,a4paper,oneside]{extbook}
\usepackage[utf8]{vietnam}
\usepackage{amsmath,amssymb,mathrsfs,fancyhdr,enumerate,multirow,makecell}
\usepackage{tkz-euclide,tikz-3dplot,tikz,tkz-tab}
\usetikzlibrary{arrows.meta,calc,intersections,angles,quotes,backgrounds,shapes.geometric,patterns,through,shadows,positioning}
\usepackage{yhmath}
\usepackage[top=0.7cm, bottom=0.7cm, left=1.0cm, right=1.0cm, headheight=16pt, headsep=3mm, footskip=5mm]{geometry}
\usepackage[hidelinks,unicode]{hyperref}
\usepackage{qrcode}
\usepackage[many]{tcolorbox}
\usepackage{fontawesome}

\pagestyle{fancy}
\fancyhf{}
\fancyhead[L]{\footnotesize\bfseries\color{teal!85!black}\faBook~SỔ TAY KIẾN THỨC TOÁN 4}
\fancyhead[R]{\footnotesize\bfseries\color{blue!85!black}Lớp Toán Thầy Lĩnh -- loptoanthaylinh.vn}
\fancyfoot[C]{\footnotesize\bfseries\color{gray!80!black}-- Sổ tay kiến thức: Bài 51. Số lần xuất hiện của một sự kiện --}
\renewcommand{\headrulewidth}{0.4pt}

\begin{document}

% --- BANNER TIÊU ĐỀ BÀI HỌC ---
\noindent
\begin{tcolorbox}[
  enhanced,
  colback=teal!8!white,
  colframe=teal!80!black,
  arc=3mm, boxrule=1pt,
  left=3mm, right=3mm, top=2mm, bottom=2mm,
  halign=center
]
  {\Large\bfseries\color{teal!90!black}\faBookmark\ BÀI 51: SỐ LẦN XUẤT HIỆN CỦA MỘT SỰ KIỆN}\\[1.5pt]
  {\footnotesize\itshape\color{black!75} Thực nghiệm kiểm đếm, ghi chép tần số xuất hiện và nhận xét khả năng xảy ra trên 1 trang A4}
\end{tcolorbox}

\vspace{-2.5mm}

\noindent
\begin{minipage}[t]{0.49\textwidth}\vspace{0pt}
\begin{tcolorbox}[
  enhanced,
  colback=blue!4!white,
  colframe=blue!75!black,
  arc=2mm, boxrule=0.8pt,
  left=2mm, right=2mm, top=1.8mm, bottom=1.8mm,
  title={\footnotesize\bfseries\color{white}\faRepeat\ 1. Kiểm đếm thực nghiệm lặp lại}
]
{\footnotesize
Khi thực hiện lặp lại một hoạt động nhiều lần (tung đồng xu, gieo xúc xắc, lấy bóng trong hộp kín):
\begin{itemize}\setlength{\itemsep}{1pt}\setlength{\leftskip}{-2.5mm}
  \item Mỗi kết quả mong muốn xảy ra gọi là một \textbf{sự kiện}.
  \item Mỗi lần sự kiện xảy ra, ta ghi lại bằng \textbf{1 vạch} ($|$).
  \item Quy ước đếm tiện lợi: Cứ $5$ lần thì gạch chéo cụm $4$ vạch thành một bó:
  $$\text{Cụm 5 vạch: } \rlap{\Big/}|~|~|~|$$
  \item Tổng số vạch chính là \textbf{số lần xuất hiện} của sự kiện đó.
\end{itemize}
}
\end{tcolorbox}
\end{minipage}
\hfill
\begin{minipage}[t]{0.49\textwidth}\vspace{0pt}
\begin{tcolorbox}[
  enhanced,
  colback=teal!4!white,
  colframe=teal!75!black,
  arc=2mm, boxrule=0.8pt,
  left=2mm, right=2mm, top=1.8mm, bottom=1.8mm,
  title={\footnotesize\bfseries\color{white}\faCheckSquareO\ 2. Nhận xét khả năng xảy ra}
]
{\footnotesize
\begin{itemize}\setlength{\itemsep}{1pt}\setlength{\leftskip}{-2.5mm}
  \item \textbf{Định luật bảo toàn số lần:}
  $$\text{Tổng số lần xuất hiện} = \text{Tổng số lần thử nghiệm}$$
  \item \textbf{So sánh tần số:}
  \begin{itemize}\setlength{\itemsep}{0.5pt}\setlength{\leftskip}{-2.5mm}
    \item Sự kiện xuất hiện \textbf{nhiều lần nhất} $\leftrightarrow$ Khả năng xảy ra cao nhất trong thực nghiệm đó.
    \item Sự kiện xuất hiện \textbf{ít lần nhất} $\leftrightarrow$ Khả năng xảy ra thấp nhất.
  \end{itemize}
  \item Số lần thực nghiệm càng nhiều thì kết quả thống kê càng khách quan và chính xác.
\end{itemize}
}
\end{tcolorbox}
\end{minipage}

\vspace{-2.5mm}
% --- KHỐI 3: SƠ ĐỒ TƯ DUY TIKZ ---
\noindent
\begin{tcolorbox}[
  enhanced,
  colback=orange!4!white,
  colframe=orange!80!black,
  arc=2mm, boxrule=0.8pt,
  left=2mm, right=2mm, top=1.5mm, bottom=1.5mm,
  title={\footnotesize\bfseries\color{white}\faSitemap\ SƠ ĐỒ TƯ DUY: TIẾN TRÌNH THỰC NGHIỆM XÁC SUẤT 3 GIAI ĐOẠN}
]
\centering
\begin{tikzpicture}[
  box/.style={draw, rounded corners=2pt, align=center, font=\scriptsize\bfseries, inner sep=3pt, minimum height=1.1cm},
  arr/.style={-{Stealth[length=2.5mm]}, thick, shorten >=1.5pt, shorten <=1.5pt}
]
  \node[box, fill=blue!10, draw=blue!80!black, text width=4.3cm] (b1) at (0,0) {BƯỚC 1: THỰC NGHIỆM\\Lặp lại hoạt động $N$ lần\\(Tung xu, gieo xúc xắc)};

  \node[box, fill=teal!10, draw=teal!80!black, text width=4.6cm] (b2) at (5.4,0) {BƯỚC 2: KIỂM ĐẾM GHI CHÉP\\Dùng vạch gạch kiểm đếm\\Nhóm cụm 5 vạch};

  \node[box, fill=orange!12, draw=orange!80!black, text width=4.5cm] (b3) at (10.9,0) {BƯỚC 3: KẾT LUẬN\\Tổng hợp số lần xuất hiện\\Đánh giá Min / Max};

  \draw[arr, blue!80!black] (b1) -- (b2);
  \draw[arr, teal!80!black] (b2) -- (b3);

  \node[above, font=\tiny\bfseries\color{blue!80!black}] at (2.7, 0.1) {Mỗi lần xảy ra};
  \node[above, font=\tiny\bfseries\color{teal!80!black}] at (8.15, 0.1) {Cộng dồn vạch};
\end{tikzpicture}
\end{tcolorbox}

\vspace{-2.5mm}
% --- KHỐI 4: VÍ DỤ & MẸO TRÁNH BẪY ---
\noindent
\begin{minipage}[t]{0.49\textwidth}\vspace{0pt}
\begin{tcolorbox}[
  enhanced,
  colback=green!4!white,
  colframe=green!60!black,
  arc=2mm, boxrule=0.8pt,
  left=2mm, right=2mm, top=1.8mm, bottom=1.8mm,
  title={\footnotesize\bfseries\color{white}\faCheckCircle\ Ví dụ thực hành chuẩn mực}
]
{\footnotesize
\textbf{Ví dụ:} Bạn Nam tung đồng xu 20 lần và ghi lại:
\begin{itemize}\setlength{\itemsep}{1pt}\setlength{\leftskip}{-2.5mm}
  \item Mặt Sấp (S): xuất hiện 12 lần.
  \item Mặt Ngửa (N): xuất hiện 8 lần.
\end{itemize}
\textbf{Nhận xét kết quả:}
\begin{itemize}\setlength{\itemsep}{1pt}\setlength{\leftskip}{-2.5mm}
  \item Mặt Sấp xuất hiện \textbf{nhiều hơn} mặt Ngửa: $12 - 8 = \mathbf{4}$ lần.
  \item Tổng số lần xuất hiện: $12 + 8 = \mathbf{20}$ lần (bằng đúng số lần tung ban đầu).
\end{itemize}
}
\end{tcolorbox}
\end{minipage}
\hfill
\begin{minipage}[t]{0.49\textwidth}\vspace{0pt}
\begin{tcolorbox}[
  enhanced,
  colback=red!4!white,
  colframe=red!70!black,
  arc=2mm, boxrule=0.8pt,
  left=2mm, right=2mm, top=1.8mm, bottom=1.8mm,
  title={\footnotesize\bfseries\color{white}\faExclamationTriangle\ Mẹo tránh bẫy trắc nghiệm}
]
{\footnotesize
\begin{itemize}\setlength{\itemsep}{1.5pt}\setlength{\leftskip}{-2.5mm}
  \item \textbf{Bẫy đếm thiếu số lần thực nghiệm:}\\
  Sau khi kiểm đếm, luôn cộng tổng các sự kiện lại xem có bằng tổng số lần thử nghiệm đề bài cho không. Nếu không bằng $\to$ Đếm sót hoặc đếm thừa!
  \item \textbf{Bẫy nhầm giữa "Mặt" và "Số lần":}\\
  Ví dụ gieo xúc xắc được mặt 6 chấm 2 lần, học sinh hay nhầm lấy số 6 thay vì số lần xuất hiện là 2.
\end{itemize}
}
\end{tcolorbox}
\end{minipage}

\end{document}
